\clearpage
\properlane{gift-becomes-charity}{The invitation becomes a life of charity}
The feast is God's gift, but it is a gift of communion. To accept it is to be drawn into the love of the King and his Son. Entry cannot be earned by an impressive past; neither can a person receive communion as though it imposed no claim on the present. Gregory's garment interpretation brings this tension into focus. Augustine's teaching on forgiveness and purification shows how mercy can be both the beginning of the Christian life and the power of its continuing change.

\subsection{Why the wedding is the present Church}
Gregory starts with a feature easily passed over: a guest is expelled. If the wedding simply represented the completed heavenly assembly, his expulsion would be difficult to explain. The final feast admits no subsequent loss of its happiness. The banquet in Matthew therefore signifies, in Gregory's principal reading, the present Church, where bad and good are gathered and where a person can enter without persevering. He contrasts this with the ultimate banquet in Luke, while allowing that the two accounts might also be reconciled without a dispute about faith.\footnote{Gregory, \work{Homiliae in Evangelia} 38.1, 7.} His account is an inference from the parable's movement, not merely an equation of a symbol with a familiar institution.

This location changes whom the parable addresses. The bad guest is not evidence that the invitation should never have reached him. The servants were told to gather everyone they found. Nor does the mixed company prove that the good should leave to establish a feast of their own. Gregory turns the presence of troublesome people into a test of the hearer's charity. While good and evil remain intermingled, the refusal to endure another's faults can reveal the would-be righteous person's own defect. His examples of grain amid chaff and roses amid thorns describe the difficulty of life together.\footnote{Gregory, \work{Homiliae} 38.7--8.}

The moral consequence is precise. Another person's failure does not discharge my obligation to love; my correct identification of that failure does not itself make me good. Charity can oppose wrongdoing while refusing hatred of the person. The wedding's openness and its judgment belong together because both concern communion: God gathers persons who need transformation, and their refusal of communion's life matters. A community that treats moral failure as harmless loses the garment; a community that refuses all imperfect guests forgets how the room was filled.

The short Gospel preserves the gathered bad and good but stops before the inspection. Gregory's argument from expulsion is consequently an explanation of the \emph{long} Gospel. When the short form is proclaimed, the immediate emphasis falls on responding to the call and sharing the room with unexpected people. Charity can still be developed through the other appointments and the Gospel's wider context, but the absent garment scene cannot be treated as that proclamation's concluding image.

\subsection{Why baptism or faith alone is not the garment}
Gregory's next question is what the guest lacks. He tests two answers: baptism and faith. In his ecclesial interpretation these already explain the person's entry. The garment must therefore name something that can be absent in someone who has entered. He identifies it as charity. The point is not that baptism is empty or faith unnecessary; the point is that a person who has been admitted can still fail to live the communion received.\footnote{Gregory, \work{Homiliae} 38.9.}

Gregory then gives charity a Christological reason. Love brought the Son to unite human beings to himself. It is fitting that the same love characterize those who celebrate his wedding. The garment is not an arbitrary condition added by a capricious host after the guests arrive. Within this interpretation it expresses the very reality for which the feast is held. The Son joins the Church to himself; a heart closed to God and neighbor contradicts that union from within.

Chrysostom reaches a compatible result by distinguishing the calling from the life that follows it. Calling and cleansing are gifts of grace; the recipient must take care over the ensuing life. He names the garment as life and practice. The new guests have no warrant to make faith into an exemption from the judgment of their actions. His questions turn the parable against the complacency of those who might think the earlier guests' rejection has made them automatically secure.\footnote{Chrysostom, \work{Homilies on Matthew} LXIX.2.} Gregory identifies the inner form of a fitting life; Chrysostom stresses its visible practice. They do not describe two competing currencies with which to purchase admission.

The distinction also guards the meaning of the man's silence. Gregory understands it as the collapse of excuse before one who knows the conscience. Chrysostom sees the absence of an answer as the guest's own condemnation of himself. Neither needs an invented story about a royal cloakroom. The garment image discloses a moral relation to the feast, not an undocumented clothing custom. Its force lies in the discrepancy between receiving an honor and continuing a life at odds with it.

\subsection{Mercy is the beginning, not the cancellation, of change}
The Entrance gives this demanding Gospel an answering word. Augustine's sinner has no secure platform of innocence from which to approach God. Hope rests in forgiveness and Christ's redemption. Yet the law of love in his psalm exposition both forgives the past and warns concerning the future; it becomes a companion on the way.\footnote{Augustine, \work{Enarrationes} 129.2--3.} Mercy has a direction. It does not simply move a person from an unfavorable legal column to a favorable one and leave the rest of life untouched.

The Collect's concern with grace before and after action belongs here. The necessary good is neither an independent achievement presented to God nor an unreal activity which God does while the person remains inert. Augustine explains the next verse after the Johannine Communion in exactly this double way. The believer purifies himself, but not by himself: the divine indwelling enables a willing response. Asking for help presumes both need and an act in which help is received.\footnote{Augustine, \work{First Epistle of John} IV.7, on 1 Jn 3:3. The verse is context, not part of the Communion excerpt.}

The psalm's Shepherd likewise leads in right paths and nourishes a growing flock. Augustine's sacramental reading connects regeneration with discipline and a mature table; Thomas joins inward conversion to outward good works.\footnote{Augustine, \work{Enarrationes} 22.2--5; Thomas, \work{Super Psalmo} 22, n. 1.} The direction of the path cannot be detached from the food offered along it. These elements make the garment's demand intelligible as a life sustained by gifts, rather than a demand to arrive already perfect before receiving any.

Isaiah's all-peoples feast widens this life beyond a circle of familiar guests. Thomas's link to Christ's Passion places its abundance under the saving action from which the guests receive everything. The end of death and tears also prevents the present mixed condition from becoming a permanent definition of God's people. Charity is practiced amid imperfection because the Lord is forming a communion whose finished life will not be wounded by death.

\subsection{The twofold weave: God and the neighbor}
Gregory compares charity to a garment woven on two beams, love of God and love of neighbor. Neither can replace the other. Contemplation cannot become an excuse to abandon a person who needs help; service cannot become a pretext for losing the God toward whom the journey tends. His image of twice-dyed scarlet gives the same warning: the one flame must burn in both directions.\footnote{Gregory, \work{Homiliae} 38.10.} This is not a schedule allocating some hours to God and others to people. It describes the unity of the love with which both are sought.

The test becomes sharper in Gregory's account of enemies. Affection for someone who loves us can be mere reciprocity. Charity loves the friend in God and the enemy for God's sake. Hatred, envy of another's good and the intention secretly to injure someone disclose a torn garment even in a person accustomed to the feast.\footnote{Gregory, \work{Homiliae} 38.11.} The acclamation's prayer for enlightened hearts thus reaches beyond being able to explain a verse. A person may understand the external order of the celebration while failing to see what communion asks of a concealed resentment.

Philippians gives charity bodily form. The donors share Paul's distress; the recipient recognizes their gift without letting either self-sufficiency or patronage govern the relation. When the assembly brings offerings, those gifts can be understood within the same movement toward fellowship and heavenly glory. Liturgical giving does not make ordinary generosity unnecessary. Nor can an offering purchase the right to continue hostility toward another guest. The good sought in the psalm Communion includes, in Augustine's explanation, charity itself among the riches of the heart.\footnote{Augustine, \work{Enarrationes} 33, second exposition, 14.}

Gregory gives a particularly searching account of the bound hands and feet. The feet that decline to visit the sick and the hands that give nothing to those in need are already bound by a person's own refusal. Judgment makes manifest the bondage chosen in life.\footnote{Gregory, \work{Homiliae} 38.13.} The image refuses an exclusively interior account of love. A feeling of goodwill does not yet visit anyone. The garment becomes visible in actions possible to the body: going, giving, helping and refusing the use of those same powers for harm.

\subsection{Imperfect charity and the hope of sight}
The severity of this account raises a further difficulty. If the garment is charity, must someone whose love remains imperfect despair? Gregory explicitly says no. He distinguishes having the garment imperfectly from lacking it altogether and directs the imperfect toward hope in the merciful King. He then warns against presumption about perseverance.\footnote{Gregory, \work{Homiliae} 38.12, 14.} The two statements serve different people and different temptations. Fear should awaken a settled refusal of love; it should not persuade a struggling person that growth has become impossible.

Augustine's account of John's promised vision supplies the positive movement of that growth. Divine childhood is real now, yet its full manifestation remains ahead. The heart's desire can enlarge its capacity; purification removes what resists the gift. The future is therefore neither merely a reward externally appended to an unchanged life nor a possession one may claim to have completed already.\footnote{Augustine, \work{First Epistle of John} IV.5--7.} With the Johannine Communion option, this relation between received identity and awaited likeness is explicit. With the psalm option, the enduring good of seeking God carries the same life forward under another emphasis.

The Prayer after Communion places nourishment in relation to divine participation. Received life is meant to become lived love and, finally, unveiled communion. The long Gospel's judgment and John's promised sight are not rivals, one frightening and the other comforting. Both make the heart's relation to God consequential. The invitation is generous enough to reach the undeserving and holy enough to change them.

\pagebreak
\begin{fourSenses}
\item[Literal.] God prepares and calls; people refuse, answer and gather. The long Gospel alone proceeds to inspect a guest. The other appointments name pardon, right paths, shared distress, illumination and a future likeness not yet manifested.
\item[Allegorical.] Gregory's present Church is gathered for union with the Son. Augustine's Christ redeems and shepherds it. Charity is the life fitting to that union, sustained by grace and sacramental nourishment.
\item[Moral.] Mercy becomes concrete love of God and neighbor, reaching the enemy, the sick and the inconvenient guest. Neither religious presence nor generosity with possessions excuses a refusal of charity. Imperfect love is called to growth, not despair.
\item[Anagogical.] The mixed gathering moves toward final judgment and perfected communion. Perseverance and purified desire prepare for divine sight; likeness fulfills adoption without making the creature equal to God.
\end{fourSenses}
